\chapter{Simulation}
\label{ch:sim}

\pkg{dharma_sim} replaces the robot and the LiDAR-inertial estimator so the whole stack runs
without hardware, in CI and on a laptop. It is being refactored in parallel with this edition;
this chapter follows its design (task file) and is checked against the finished code in phase 2.
\todo{verify against dharma\_sim 0.2.0: nodes, parameters, scenario format, generated parameter
table}

\section{Why the simulator models errors}

The 0.1 simulator had perfect odometry and stamped everything with \code{now()} (finding S1). The
code paths that handle drift (localization corrections, L1), map smearing (M1), a tilted start
(O1) and stale TF (T1, T3) were therefore never exercised. Version 0.2.0 models them, with a seed
so every run is reproducible.

\section{Models}

\begin{description}
  \item[Body.] Omnidirectional or differential from \param{robot.kinematics}, with the acceleration
    limits of the profile. It integrates the measured time step (clamped; 0.1 used a fixed step per
    wall-timer tick, so the simulated robot slowed down under load, review \#14), stops on stale
    commands and detects collisions: \topic{sim/collision} (Bool, latched true on contact) and the
    body stops.
  \item[World.] Walls, pillars (circles) and timed dynamic obstacles; ray casting that is correct
    for origins inside circles (0.1 used $t = -b - \sqrt{\mathrm{disc}}$ there and was blind,
    review \#10).
  \item[LiDAR.] A Mid-360 model: vertical field of view from $-7^\circ$ to $52^\circ$ at the mount
    height, minimum and maximum range, noise, dropout windows; a dense enough cloud (perception
    needs at least \param{blind.min_points} points per scan).
  \item[Odometry drift.] $\T{odom}{base} \leftarrow \T{odom}{base}\exp(\vect{\xi})$,
    $\vect{\xi} \sim \mathcal{N}(\vect{0}, \boldsymbol{\Sigma}\,\Delta s)$ per distance travelled,
    plus a yaw bias, a scale error and slip; seeded RNG.
  \item[Latency.] Stamps $\tau = t - L$ for LiDAR and odometry, and a command latency queue.
  \item[IMU mounting.] The Livox manual gives the IMU at $(0.011, 0.02329, -0.04412)$~m in
    \frm{livox_frame}; 0.1 used FAST-LIO's $T_{IL}$ with the opposite sign (review \#11).
\end{description}

\section{Robot-side contract}

The simulator publishes what a robot and its estimator publish: \topic{/<ns>/odom},
\topic{/<ns>/battery}, \topic{/Odometry}, \topic{/cloud_registered_body}, TF
\frm{camera_init}$\to$\frm{body} and \frm{base_link}$\to$\frm{livox_frame}$\to$\frm{livox_imu_frame},
and subscribes \topic{/<ns>/cmd_vel}. It also publishes \topic{sim/ground_truth},
\topic{sim/obstacles} and \topic{sim/inspect}, and answers \topic{/initialpose} while
\topic{localization/status} is \code{NONE}. Odometry runs with \code{method:=external}. Perception
publishes \code{BLOCKED} scans unless it gets \frm{odom}$\to$\frm{base_link} through the bridge
chain; the mux forwards nothing without a fresh, stamped, 360° safety scan in \frm{base_link}.

\section{Scenarios and site fixture}

Scenarios are YAML files with a seed, timed obstacles, dropouts and latencies, loaded by the node
and by tests (\code{config/scenarios/*.yaml}; worlds in \code{config/worlds/}, routes in
\code{config/routes/}). A site fixture writes a loadable site (map, tags and a route) from a
world, without points of interest or tours.

\section{Evaluation protocol}
\label{sec:eval-protocol}

The experiments of the paper run on these scenarios: E1 stopping distance against $D(v)$, E2
blind LiDAR in AUTO and REMOTE, E3 localization error with and without the EKF and scan matching
under seeded drift, E4 path shift clearance and completion time with standing and walking people.
\todo{scenario files, seeds, runs per condition, metrics scripts and results (phase 2)}
